
\subsection{Uncertainty quantification}

A point estimate is not enough for maintenance scheduling, so we also measure
whether the ensemble knows when it is wrong. Two interval constructions are
compared: raw ensemble sigma, $\hat y\pm z\hat\sigma$ with $\hat\sigma$ the
sample standard deviation over the fifteen models, and split conformal
calibration (CQR) whose half-width is the $(1-\alpha)$ quantile of absolute
residuals on the held-out folds.

Raw ensemble intervals under-cover badly: a nominal 90\% interval covers
61--74\% of engines, and a nominal 95\% interval covers 69--80\%. Fifteen
models that share an architecture, a dataset, and a training protocol make
highly correlated errors, so their spread systematically understates the true
error. CQR restores coverage to 0.82--0.88 at the 90\% level and 0.90--0.94 at
95\%, but the intervals are 25--31 cycles wide at 90\% and 34--45 at 95\%,
against a target range of 0--125. The cause is the heavy tail already visible
in the bias analysis: the largest single error exceeds 60 cycles.

The honest reading is that this model is strong on point accuracy and weak on
uncertainty. A 25-cycle interval does not support a maintenance decision. We
report the numbers rather than a calibration trick that hides them.
